\documentclass[11pt]{article}
\usepackage[margin=1in]{geometry}
\usepackage{amsmath,amssymb,amsfonts,mathtools}
\usepackage{booktabs,array,longtable}
\usepackage{lmodern}
\usepackage[T1]{fontenc}
\usepackage{hyperref}
\usepackage{xcolor}
\usepackage{enumitem}
\hypersetup{colorlinks=true,linkcolor=blue!50!black,urlcolor=blue!50!black}
\setlength{\parskip}{.5em}\setlength{\parindent}{0pt}
\DeclareMathOperator{\softmin}{softmin}
\DeclareMathOperator{\clamp}{clamp}
\DeclareMathOperator{\capone}{cap}
\newcommand{\EV}[1]{\mathbb{E}\!\left[#1\right]}
\newcommand{\fit}{\textbf{[fitted]}}
\newcommand{\agg}{\textbf{[aggregate]}}
\newcommand{\jud}{\textbf{[judgment]}}
\title{\textbf{The Omni Method}\\[2pt]\large Model Specification and Mathematics}
\author{Generated from the console source (\texttt{index.html})}
\date{November 21, 2026}
\begin{document}
\maketitle
\tableofcontents
\newpage

\section{Purpose and provenance labels}
The Omni Method is a valuation and saleability console for business brokers. From a business's industry, revenue, seller's discretionary earnings (SDE), margins and operating characteristics it produces: a value and its range; the probability a sale closes; time to close; post-sale survival risk (SBA loans); a buyer universe; an SBA financing check; the value gap and the interventions that pay; and the optimizer's asks.

Every parameter carries one of three labels, used throughout this document:
\begin{itemize}[nosep]
\item \fit{} --- a regression or hazard model was run on data.
\item \agg{} --- a published aggregate statistic was read off (Pepperdine, IBBA, SBA).
\item \jud{} --- a person chose it; it structures the output but is not estimated.
\end{itemize}
Where a figure rests on judgment the console says so. Notation: $E$ is SDE (``ern''), $R$ revenue, $\text{ask}$ the asking price, $V=\EV{V\mid X}$ the model value, $\Phi$ the standard normal CDF.

\section{Inputs}
Industry (one of 16: home services, medical, veterinary, MSP, SaaS, professional services, manufacturing, distribution, e-commerce, staffing, landscaping, car wash, auto, construction, fitness, restaurant); revenue $R$; SDE $E$; asking price; and nine operating dimensions: recurring revenue share ($rec$), gross margin ($gm$), customer concentration ($cc$, top-customer share), growth ($grw$), retention ($ret$), owner dependence ($own$), management depth ($mgt$), financial-record quality ($fin$), channel concentration ($plt$); plus years of clean history, number of buyers contacted $n$, and the seller's cost-of-delay $\lambda_T$ and fallback value $\theta$.

\section{Peer standardization}
Each dimension is compared with the industry peer median $\bar x_{\text{peer}}$ and a peer standard deviation $\sigma_{\text{peer}}$:
\[
z_i=\frac{x_i-\bar x_{\text{peer}}}{\sigma_{\text{peer}}},\qquad z_i\leftarrow\clamp(z_i,-3.2,3.2),\qquad \text{percentile}=\Phi(z_i).
\]
Peer standard deviations \fit/\agg: $rec\,0.28$, $cc\,0.18$, $own\,0.25$, $mgt\,0.28$, $grw\,0.16$, $gm\,0.13$, $ret\,0.13$, $fin\,0.28$, $plt\,0.22$; earnings scale $0.48$ (in $\ln$).

\section{Value model (Model II)}
\subsection{Fitted regression}
Fitted to $n=18{,}941$ closed sales (BIZCOMPS); target $\ln(\text{price in \$000})$; cross-validated $R^2$ on log price $0.755$; median absolute percentage error $28.4\%$; $44.5\%$ within $\pm25\%$. \fit
\begin{align}
\ln\hat P_{\$000}&=0.7637+0.3117\ln\tfrac{E}{\$1\text{k}}+0.4057\,h(E)+0.4709\ln\tfrac{R}{\$1\text{k}}\notag\\
&\quad+0.4097\,m+0.005516\,(y-2015)+1.3220\,e_i,\qquad m=\clamp(E/R,0,1),\;y=2026.
\end{align}
The size hinge, with knot $\$150$k and cap $\$3$M of SDE (in thousands, $k_0=150$, $k_c=3000$):
\begin{equation}
h(E)=\begin{cases}\max\!\big(\ln\tfrac{E_k}{k_0},0\big) & E_k\le k_c\\[2pt]
\ln\tfrac{k_c}{k_0}+0.5360\,\ln\tfrac{E_k}{k_c} & E_k>k_c\end{cases}
\end{equation}
Above the cap the slope $0.536$ makes the total scale elasticity exactly one, so the implied multiple goes flat rather than climbing or falling off the data (only 56 of 18,941 rows sit beyond the cap). $e_i$ is the industry effect: a 6-digit NAICS effect shrunk toward its 2-digit sector, falling back to the sector effect and then the global effect.

\subsection{Median calibration}
OLS forces the \emph{mean} log residual to zero, not the median, which ran from $-0.057$ at a \$140k prediction to $+0.101$ at \$320k. A correction $c(\cdot)$, piecewise linear in the uncorrected $\ln\hat P$ through 15 knots (from $-0.027$ at the low end to $+0.108$ at the peak, then $+0.064$ at the top), is added:
\begin{equation}
\ln\hat P^{c}=\ln\hat P+c(\ln\hat P),\qquad \hat P^{c}\text{ then has median ratio }1.000.
\end{equation}
\subsection{Quality position inside the measured spread}
The regression sees earnings, revenue and industry but not recurring revenue or owner dependence. What it does measure is the scatter of sales around its prediction: out-of-fold quantiles of $P/\hat P^c$ (after calibration, $n=18{,}941$) are
\begin{center}\small
\begin{tabular}{lcccccccccc}\toprule
percentile & 2 & 5 & 10 & 25 & 50 & 75 & 90 & 95 & 98\\\midrule
ratio & 0.276 & 0.397 & 0.510 & 0.727 & 1.000 & 1.318 & 1.694 & 2.046 & 2.738\\\bottomrule
\end{tabular}\end{center}
Quality decides where in this band a business sits. Composite quality $z$ is a weighted sum over nine dimensions, normalized by the weights' norm so it stays standard normal:
\begin{equation}
q_z=\frac{\sum_k w_k\,\sigma_k z_k}{\sqrt{\sum_k w_k^2}},\qquad \sigma_k=-1\text{ for }own,cc,plt;\;+1\text{ otherwise},\qquad q_p=\Phi\big(\clamp(q_z,-3,3)\big).
\end{equation}
Weights: $w_{own}=1.24$ \fit, $w_{mgt}=0.60$ \fit\ (weakly supported), all others $1$ \jud; $\sqrt{\sum w_k^2}=2.983$. The position factor, piecewise log-linear so the midpoint is exactly $1$:
\begin{equation}
q=\begin{cases}\exp\!\big[\ln(0.7273)\,(1-2q_p)\big] & q_p\le\tfrac12\\ \exp\!\big[\ln(1.3181)\,(2q_p-1)\big] & q_p>\tfrac12\end{cases}
\end{equation}
\begin{equation}
\boxed{\;V=\EV{V\mid X}=\$1\text{k}\cdot\hat P^{c}_{\$000}\cdot q\;},\qquad M=V/E.
\end{equation}
The quality position (which weights each dimension, and the placement itself) is \jud; the spread it is placed in is data.

\subsection{Range (the fan)}
Two different spreads are used, and they must not be confused.

\textbf{The reported value range} is the calibrated fitted point times the \emph{full} measured out-of-fold spread, with no quality position applied:
\begin{equation}
\begin{gathered}p_q=\hat P^{c}\cdot R(q),\qquad q\in\{10,25,75,90\}\\ R(10)=0.5105,\ R(25)=0.7277,\ R(75)=1.3188,\ R(90)=1.6949\end{gathered}
\end{equation}
So the range is centred on the pre-quality point $\hat P^c$; the model value $V=q\hat P^c$ sits inside it, shifted by the quality factor. The ratio of mean to median for this full spread is $1.0907$.

\textbf{The price distribution given a close} (used for $\EV{P\mid\text{close}}$, the price-if-it-closes quantiles and expected proceeds) keeps only part of the spread once the business is listed and priced: $W=V\cdot R(p)^{s}$ with $s=0.8$ \jud. It is carried on 48 equal-probability slices, each at the slice's exact mean (for $\ln R$ linear in $z=\Phi^{-1}(p)$ on a piece, $\mathbb E[R;a<z<b]=\exp(c+B^2/2)\,[\Phi(b-B)-\Phi(a-B)]$), then passed through the ask-anchoring map $P(W)$ of Section~5.

\subsection{Confidence and extrapolation}
The sliders reach beyond the data (SDE to \$20M against a 99th percentile of \$1.7M). The console reports a confidence label from the sampling error of the fit and flags when an input lies more than $3\times$ outside the fitted range; extrapolation is a label, not a different formula.

\section{Asking-price anchoring and the price if it closes}
Let $\rho=\text{ask}/V$, $\tilde\rho=\softmin(\max(\rho,0.35),1.5)$, $\bar a=\softmin(\text{ask},1.5\,V)$.
\begin{equation}
A(\rho)=1+(A_{\max}-1)\tanh\!\Big(\frac{\alpha\ln\tilde\rho}{A_{\max}-1}\Big),\quad A_{\max}=1.3,\ \alpha\ \jud
\end{equation}
\begin{equation}
P(W)=\softmin\!\Big(W\cdot A\!\big(\tfrac{\bar a}{W}\big)\cdot(1+\text{comp}),\ \bar a\,c_{\text{ask}}\Big),\qquad
\EV{P\mid\text{close}}=\tfrac{1}{48}\sum_{p}P(W_p).
\end{equation}
$A$ is the anchoring effect of the ask on the closing price; $c_{\text{ask}}$ is the ceiling on price relative to the ask. The soft minimum, used everywhere a hard $\min$ would make a slider feel broken, is
\begin{equation}
\softmin(x,y)=\exp\!\Big(\min(a,b)-K\ln\!\big(1+e^{-|a-b|/K}\big)\Big),\ a=\ln x,\ b=\ln y,\ K=0.012.
\end{equation}
\textbf{Competition premium} (\jud\ scale $\kappa=0.03$, $H=1.6235$):
\begin{equation}
\text{comp}=\kappa\ln\frac{H+n_{\text{offers}}}{H+n_{\text{typical}}}\in[-0.08,0.14],\qquad n_{\text{offers}}=n\cdot\overline{P_E}\cdot s_E,
\end{equation}
with $s_E=0.184$ (engagement-to-offers share, \fit; 95\% CI $0.13$--$0.226$) and $\overline{P_E}$ the buyers' mean engagement probability.

\section{Probability the sale closes}
\subsection{Reported model (aggregate)}
Pepperdine publishes termination \emph{rates}, not deal rows, so the reported probability is a base rate times a price-gap penalty times a thin-pool factor times a bounded quality factor. \agg\ + \jud
\begin{align}
P_C&=\min\!\big(0.95,\ b(V_{\text{base}})\,(1-\pi(g))\,\tau(n)\,\kappa(q_z)\big),\qquad g=100\cdot\frac{\text{ask}-V}{V}\ (\%)\\
b&=b_B+(b_I-b_B)\,u,\quad u=\clamp\!\Big(\frac{\ln(V_{\text{base}}/\$1\text{M})}{\ln 5},0,1\Big),\ b_B=0.608,\ b_I=0.6967\notag
\end{align}
where $b_B,b_I$ are the recent-three-observation means of the broker-listed and intermediated close rates (pooled 2011--2024 values are $0.6128$ and $0.708$; intermediated termination troughed at 24\% in 2020--21 and is 32\% in 2026, so the recent mean is used). Channel overrides: $u=1$ if the deal is intermediated, $0$ if a broker listing.

\textbf{Price-gap penalty.} Pepperdine publishes the distribution of gaps at which price-killed deals died. That is the spread of gaps \emph{among failed deals}, not the probability that a deal fails at a given gap (most failures sit at 11--30\% largely because many asking prices do too), so it does not identify a failure rate. The console borrows its cumulative \emph{shape} $F_{\text{gap}}$ for a monotone penalty and sets its size, $\pi_{\max}$, by judgment; the Data tab says the penalty is ``not identified''. Knots of the shape:
\begin{center}\small\begin{tabular}{lccccccc}\toprule
$g$ (\%) & 0 & 11.11 & 25 & 42.86 & 66.67 & 100 & 200\\\midrule
$F_{\text{gap}}$ & 0 & 0.0699 & 0.4276 & 0.7920 & 0.8869 & 0.9301 & 1\\\bottomrule\end{tabular}\end{center}
\begin{equation}
\pi(g)=\begin{cases}0 & g\le0\\ \pi_{\max}\,F_{\text{gap}}(g) & 0<g<200\\ 1-(1-\pi_{\max})\,e^{-k(g-200)} & g\ge200\end{cases}\qquad \pi_{\max}=0.45\ \jud,\ k=5.708\times10^{-4}.
\end{equation}
72.3\% of price-killed deals died at a gap of 11--30\% of the ask (12.4--42.9\% over value, converted by $p/(100-p)$); the console's ``kill zone'' marker is a descriptive reading of that statistic, not an estimated danger zone.

\textbf{Thin-pool factor.}
\begin{equation}
\tau(n)=\capone_1\!\Big(\frac{1-e^{-n_{\text{pool}}}}{1-e^{-n_{\text{typical}}}}\Big)\quad(\tau=1\text{ at a typical or deeper pool}),
\end{equation}
the scale of the thin-pool effect being \jud.

\textbf{Quality factor} (added in review round 11; \jud). Pepperdine publishes no closing rate split by business quality, so the effect is bounded by what the survey does measure: the failure reasons quality bears on. Seller misrepresentation (4.25\% of failures) and insufficient cash flow (6.0\%) together are about a tenth of all failures. A business at the best end of quality avoids that share and one at the worst end doubles it, which bounds the effect to roughly $\pm10\%$ of the base rate:
\begin{equation}
\kappa(q_z)=1+0.10\,\tanh\!\Big(\frac{\clamp(q_z,-3,3)}{1.5}\Big)\in[0.90,1.10],\qquad \kappa(0)=1,
\end{equation}
with $q_z$ the composite quality $z$ of Section~4.3 (the same weighted composite that places the value). Because $\kappa=1$ at the median business, the published base rates still describe the average deal. The effect is on top of the route quality already had: a better business has a higher $V$, hence a smaller gap $g$ at the same dollar ask. The scale $0.10$ and the width $1.5$ are judgment; the Data tab lists them as one parameter. The same factor multiplies each buyer's closing probability $P_{C,j}$ and the Deals tab's base rate when it fits the penalty scale.

\subsection{The retained feature model}
A 15-feature logistic is kept for the contribution panel but is not the reported probability:
\begin{equation}
\begin{aligned}P_C^{\text{judg}}&=\Big[p_{\min}+\tfrac{p_{\max}-p_{\min}}{1+e^{-f(X)}}\Big]\cdot G(x_{\text{price}}),\qquad f(X)=\beta_0+\sum_i\beta_ix_i,\\ G(x)&=\begin{cases}1&x\le0\\ \max(g_{\min},e^{-(g_1x+g_2x^2)})&x>0\end{cases}\end{aligned}
\end{equation}
with nonlinear terms: recurring $f_r=a\tanh(x/s)$; concentration $f_c=-(\lambda x+\kappa\max(0,x)^2)$; margin $f_m=\beta_{em}x$ for $x\ge0$ and $-(\beta_{em}|x|+0.055x^2)$ below; scale $f_z=a_z\tanh(x/s_z)$ above the median, $\beta_{size}x$ below; under-pricing $f_u=-\beta_u\min(0,x_{\text{price}})$; short history $f_y=-0.55\min\!\big(1,\max(0,\tfrac{3-\text{yrs}}{2})\big)$. All \jud; they explain, they do not drive the headline.

\section{Time to close (Model III)}
\subsection{Log-logistic duration model}
Fitted to IBBA / M\&A Source Market Pulse histograms (15 quarters, five size bands). \fit
\begin{equation}
F(t)=\frac{1}{1+(t/t_{\text{med}})^{-\beta}},\qquad S(t)=1-F(t),\qquad F(t_{\text{med}})=\tfrac12,\qquad \EV{T}=t_{\text{med}}\frac{\pi/\beta}{\sin(\pi/\beta)}.
\end{equation}
Median months from listing to close, by size band (fitted; published values in parentheses):
\begin{center}\small\begin{tabular}{lccccc}\toprule
band & $<$\$500k & \$500k--1M & \$1--2M & \$2--5M & \$5--50M\\\midrule
median (months) & 6.54 (6.5) & 7.75 (7.25) & 8.02 (8) & 9.03 (8) & 10.14 (10.5)\\
shape $\beta$ & 2.20 & 2.64 & 2.64 & 2.64 & 2.64\\\bottomrule\end{tabular}\end{center}
\subsection{Overpricing multiplier}
\begin{equation}
\begin{gathered}t_{\text{med}}=t_{\text{band}}\,e^{\eta\,h_s(x)},\qquad x=\ln\tfrac{\text{ask}}{V},\qquad \eta=0.10,\ \delta=0.10\ \ (\text{both }\jud)\\ h_s(x)=\begin{cases}0&x\le0\\ x^2/(2\delta)&0<x<\delta\\ x-\delta/2&x\ge\delta\end{cases}\end{gathered}
\end{equation}
so a listing at or below value keeps the band median, and the extra time grows smoothly (quadratic, then linear) with overpricing. These shares are conditional on the deal closing; across all listings the share closed within 365 days is about $P_C\times0.672$. The cost of delay in the optimizer is charged on $\EV{T}$ (the mean, which exceeds the median), not on $t_{\text{med}}$.

\section{Post-sale survival risk (Model V)}
SBA 7(a) change-of-ownership term loans, FY2018--FY2026 (data.sba.gov FOIA snapshot, 2026-06-30); $n=41{,}200$ loans, 721 charge-offs, 11{,}883 paid in full, 28{,}596 open. \fit\ Cross-validated AUC $0.718$, Brier $0.0201$.

\textbf{Cause-specific piecewise-exponential hazards} over months-on-book intervals with edges every 6 months to month 60, then 72, 84 and beyond; charge-off ($\lambda_1$) competes with paid-in-full ($\lambda_2$); open loans are censored. For interval $j$:
\begin{align}
u_j&=\beta_{\ln amt}\ln(\text{loan})+\beta_{fr}\,\text{franchise}+\beta_{sp}\,\text{spread}\notag\\&\quad+\beta_{\ln amt\times age}\ln(\text{loan})\Big(\tfrac{\min(\text{age}_j,72)}{12}-4\Big)+\text{sector}+\text{naics6}+\text{state}\\
\lambda_{1,j}&=\exp\!\big(a_{1,j}+\ell_0+\ell_1u_j\big),\qquad \lambda_{2,j}=\exp\!\big(a_{2,j}+\beta^{\text{pif}}\!\cdot x+\text{sector}^{\text{pif}}\big)\\
\text{CIF}(H)&=\sum_j S_j\frac{\lambda_{1,j}}{\lambda_{1,j}+\lambda_{2,j}}\Big(1-e^{-(\lambda_{1,j}+\lambda_{2,j})\Delta_j}\Big),\notag\\ S_{j+1}&=S_j\,e^{-(\lambda_{1,j}+\lambda_{2,j})\Delta_j}
\end{align}
Charge-off coefficients: $\beta_{\ln amt}=-0.2287$, $\beta_{fr}=-0.0167$, $\beta_{sp}=0.9434$, $\beta_{\ln amt\times age}=0.1586$; shrinkage $\ell_1=0.7433$, shift $\ell_0=-0.0093$. Sector effects range from $-0.389$ (health care) to $+0.485$ (transportation); NAICS-6 and state effects are gamma-shrunk (random effects). Paid-in-full: $\beta_{\ln amt}=-0.054$, $\beta_{fr}=+0.126$.

\textbf{Reported figures.} The term probability and the 5-year probability are
\begin{equation}
p_{\text{term}}=\text{CIF}(96\text{ mo observed})\cdot\text{tail}(\text{loan}),\qquad \text{term}=120\text{ mo},
\end{equation}
where $\text{tail}\approx1.08$--$1.30$ by loan size band (pooled factor $1.21\pm0.06$) extends the observed 96 months to the full term. $p_5=\text{CIF}(60)$. Relative risk $\text{rel}=p_{\text{term}}/p_{\text{base}}$; tier: $<0.75$ low, $<0.95$ below average, $<1.25$ average, $<1.75$ elevated, else high.

\textbf{Loss given charge-off.}
\begin{equation}
\text{sev}(a)=\big[1+\exp\!\big(-(0.999+0.0104(\ln L-\ln L_c)-0.3055(\text{age}/12-4))\big)\big]^{-1},
\end{equation}
\begin{equation}
\EV{\text{loss}}=\sum_j \Pr(\text{charge-off in }j)\cdot\text{sev}(\text{mid}_j)+\text{tail}\cdot\text{sev}(108).
\end{equation}
The SBA interest rate and job counts are themselves fitted (rate spread a function of $\ln$ loan with prime $7.00\%$ as of 2026-09-17; spread $2.5\%$ at the small end falling to $1.42\%$ at the large).

\section{Buyer universe}
Sixteen-industry archetype table with strategic-fit (\jud). Buyer mix by deal size (individual, strategic, financial) is interpolated in $\ln(\text{deal value})$ from IBBA 2012--2023 \agg. For buyer $j$, with multiple $M$, buyer multiple relative $m_j$, size fit $s_j\in[0,1]$, synergy share $\varphi$ and synergy pool $S_j$:
\begin{align}
S_j&=\min\!\Big(\sigma_j\,\text{synE}\,M\,s_j,\ \tfrac{0.40\,E\,M}{\varphi}\Big)\quad(\text{capitalised synergy, in dollars; so }\varphi S_j\le0.40\,E\,M)\\
V_j&=E\cdot M\,\big(1+(m_j-1)s_j\big)+\varphi S_j\\
EV_j&=P_{E,j}\,P_{C,j}\,\softmin\!\big(V_j,\ \bar a\,c_{\text{ask}}\big).
\end{align}
Here $\text{synE}=\text{strat}\cdot(0.050R+0.028R\,gm+0.015R)$ is the annual synergy pool (\jud), $\sigma_j$ the archetype's synergy intensity, $M$ the model multiple, and $\varphi=0.38$ the share of the capitalised synergy a buyer pays away (\jud); $m_j$ is the archetype's multiple relative to the all-buyer norm ($0.97$ for the individual buyer, \fit-calibrated against IBBA's deal mix). The cap on $S_j$ is written with $1/\varphi$ so that the dollars paid away, $\varphi S_j$, never exceed 40\% of value. $P_{E,j}$ and $P_{C,j}$ are the buyer's engagement and closing probabilities. ``Pursue first'' is $\arg\max_j EV_j$; ``pays most'' is $\arg\max_j V_j$ among buyers actively looking.

\section{Value gap, levers and the optimizer}
\subsection{Value gap}
\begin{equation}
\text{Value Gap}=\EV{V\mid X_{\text{potential}}}-\EV{V\mid X_{\text{current}}}.
\end{equation}
\subsection{Marginal value of a quality dimension}
Central difference, reported per one peer standard deviation:
\begin{equation}
\frac{\partial y}{\partial x_i}\approx\frac{y(x_i+h)-y(x_i-h)}{2h},\quad y\in\{V,\ P_C\,\EV{P_{\text{real}}}\},\quad\text{reported}=\frac{\partial y}{\partial x_i}\sigma_i.
\end{equation}
Ask and buyer-pool steps: $\Delta y=y(0.95\,\text{ask})-y(\text{ask})$ and $y(1.25\,n)-y(n)$.
\subsection{Intervention return}
\begin{equation}
\text{ROI}(a)=\frac{\EV{\text{proceeds}\mid X+\Delta X(a),\text{repriced}}-\EV{\text{proceeds}\mid X,\text{repriced}}-C(a)}{C(a)},\quad MV_i=\frac{\partial V}{\partial a_i}.
\end{equation}
Costs $C(a)$ and durations are \jud\ (documented per intervention; e.g.\ hiring a general manager is checked against a loaded-salary alternative). Expected proceeds $=P_C\cdot\EV{P_{\text{real}}}$.
\subsection{Budgeted optimizer}
\begin{equation}
\max_a\ \Big(P_C(a)\,\EV{P_{\text{real}}(a)}-C(a)\Big)\ \text{s.t. }C(a)\le B,\qquad \mathcal L=P_C\,\EV{P_{\text{real}}}-C-\lambda(C-B).
\end{equation}
All $2^{|\text{INTERV}|}$ combinations of interventions are searched exhaustively. The seller-utility version, used for the ``patient'' and ``sell quickly'' asks, is
\begin{equation}
U(a)=\underbrace{P_C\,\EV{P_{\text{real}}}}_{\text{sale}}+\underbrace{(1-P_C)\,\theta\,V}_{\text{business retained if it fails}}-\lambda_T\big(\EV{T}+T_{\text{exec}}\big)-C(a),
\end{equation}
with $\theta$ (fallback value fraction) and $\lambda_T$ (dollars per day of delay) stated by the user. Asks are chosen over a grid of price ratios $\rho$ capped at $1.5\,V$.

\section{SBA financing check}
Rules applied to the structure (SBA SOP 50\,10\,8.1, Appendix 15, effective October 1, 2026; SBA Information Notice 5000-881797 for FY2027 fees). \agg
\begin{center}\small\begin{tabular}{lc}\toprule
rule & value\\\midrule
maximum 7(a) loan & \$5,000,000\\
minimum DSCR: initial acquisition & 1.25\\
minimum DSCR: business expansion & 1.15\\
minimum DSCR: owner buyout & 1.25\\
minimum equity injection & 10\% of project\\
standby note counted toward equity & up to 50\% of the equity\\
amortization & 10 years\\
independent valuation required above & \$350,000\\
quality-of-earnings review at or above & \$3,000,000\\\bottomrule\end{tabular}\end{center}
With cash share $c$, seller-note share $\nu$ (default $c=10\%$, $\nu=5\%$, 6\% for 10 years), extra project costs $x$, and price $P$ (stake share $\varsigma$ for buyouts):
\begin{align}
\text{base}&=(\varsigma P+x)(1-c)-\nu\varsigma P,\qquad L=\text{base}+(1-c)\,\text{fee}(L)\ \text{(solved by iteration)},\ L\le\$5\text{M}\\
\text{project}&=\varsigma P+x+\text{fee}(L),\qquad \text{cash}=\text{project}-L-\text{note}\\
\text{DS}&=\operatorname{pmt}(L,r_{\text{SBA}},10)+\operatorname{pmt}(\text{note},r_{\text{note}},N_{\text{note}})\ \ (\text{note payment}=0\text{ on full standby})\\
\text{CF}&=E-\text{draw}-\text{capex},\qquad \text{DSCR}=\text{CF}/\text{DS}\ \ge\ \text{DSCR}_{\min}\\
\text{equity}&\ge0.10\cdot\text{project}\quad(\text{buyout: }0.10\cdot\varsigma P)
\end{align}
The owner draw is a market pay for the new owner (an estimate \jud) unless the user supplies one. The SBA guaranty fee is $0.02\cdot0.85\,L$ up to \$150k; $0.03\cdot0.75\,L$ up to \$700k; and $0.035\cdot\min(0.75L,\$1\text{M})+0.0375\cdot\max(0,0.75L-\$1\text{M})$ above. It is financed in the loan by default; it is waived for a loan of \$700k or less when the user marks the waiver or the industry is manufacturing, and the fee on a larger loan is still charged. The highest financeable price is found by solving the DSCR, equity and loan-cap conditions together; a debt above the model value is flagged.

\section{Worked example}
Inputs (the framework example): medical practice, $R=\$5{,}997{,}911$, $E=\$1{,}448{,}772$, ask $=\$3{,}221{,}069$, growth $14\%$, gross margin $64\%$, recurring $72\%$, concentration $18\%$, retention $91\%$, channel concentration $38\%$, owner dependence $72\%$, management depth $36\%$, financial-record quality $74\%$, six years of clean history, 65 buyers; the risk model uses the fitted SBA rate for this price (8.63\%: prime $7.00\%$ plus the fitted spread on 85\% of the price), and the financing check uses the rate for the actual loan size (8.89\%).

\begin{center}\small\begin{tabular}{lr}\toprule
quantity & value\\\midrule
size hinge $h$ & 2.268\\
industry effect $e_i$ & $-0.0455$\\
fitted point $\hat P^c$ (before quality) & \$3,684,729\\
median calibration & 1.066\\
quality $z$, percentile & 0.735, 76.9\%\\
quality factor $q$ & 1.160\\
\textbf{model value $V$} & \textbf{\$4,274,642} ($2.95\times$ SDE)\\
p10 / p25 / p75 / p90 (around the pre-quality point) & \$1.88M / \$2.68M / \$4.86M / \$6.25M\\
ask over value $g$ & $-24.6\%$ (penalty 0)\\
\textbf{chance of closing $P_C$} & \textbf{71.1\%} (base 68.0\%, quality $\times1.045$)\\
median time to close & 278 days (mean 356)\\
of deals that close, closed within 90 / 180 / 365 days & 4.8\% / 24.1\% / 67.2\%\\
price if it closes (median / mean) & \$3.22M / \$3.11M\\
expected proceeds $P_C\EV{P\mid\text{close}}$ & \$2.21M\\
charge-off probability to term, 5-year & 2.8\%, 0.62\% (relative risk 0.49, ``low'')\\
buyer to pursue first / pays most & search fund / direct strategic competitor\\
SBA loan (financed fee \$76,439) & \$2,806,703 at 8.89\%\\
DSCR (minimum 1.25), owner draw \$111k assumed & 3.00\\
cash required (10\% equity) & \$329,751\\\bottomrule\end{tabular}\end{center}
Check by hand: $\ln\hat P$ (before calibration) with $E_k=1448.8$, $R_k=5997.9$, $h=\ln(1448.8/150)+0=2.268$ (below the cap), $m=0.2416$:
\begin{align*}\ln\hat P&=0.7637+0.3117(7.279)+0.4057(2.268)+0.4709(8.699)\\&\quad+0.4097(0.2416)+0.005516(11)+1.3220(-0.0455)\approx 8.148\end{align*}
so $\hat P\approx e^{8.148}\approx\$3.46$M, times the median calibration ($1.066$), giving the \$3.68M point; times $q=1.160$ gives $V=\$4.27$M.

\section{What is data and what is judgment}
\begin{center}\small\begin{tabular}{p{4.3cm}p{2.2cm}p{7cm}}\toprule
output & label & basis\\\midrule
model value, range, multiple & \fit\ + \jud & BIZCOMPS regression (19k sales); the quality position inside the spread is judgment\\
chance of closing & \agg\ + \jud & Pepperdine base rates by size; the gap-penalty scale, the thin-pool factor and the $\pm10\%$ quality factor are judgment\\
time to close & \fit\ + \jud & IBBA histograms; the overpricing multiplier is judgment\\
survival risk & \fit & 41,200 SBA acquisition loans\\
quality weights & \fit\ (2 of 9) & owner dependence and management depth from broker listings; seven weights are judgment\\
business score, archetypes, best buyer & \jud & structure the conversation; magnitudes are reasoned\\
interventions and optimizer asks & \jud & costs and effects reasoned, not estimated\\
financing & \agg & SBA published rules applied to the figures; market pay is an estimate\\\bottomrule\end{tabular}\end{center}

\section{Known limits of the model}
\begin{itemize}[nosep]
\item \textbf{Typical error.} Median absolute error is 28\% on held-out sales. A gap of 20\% over value is inside that noise; the console's ask-versus-value statements are statements about the model, not the market.
\item \textbf{Quality scale.} The composite quality $z$ is standard normal only if the nine dimensions are uncorrelated. Owner dependence and management depth are correlated, so extreme positions occur more often than $\Phi$ implies.
\item \textbf{Quality moves closing odds only a little.} $P_C$ depends on size, channel, the price gap, the buyer pool and, since round 11, quality through $\kappa(q_z)$, which is bounded to $\pm10\%$ and is a judgment (no published closing rate is split by quality). Quality also acts through the value it implies. Pricing below value earns no gain in odds or time, so the optimizer's asks lean toward value or above.
\item \textbf{Value gap and intervention returns} come from the judgment quality layer, not from measurements; they are for coaching, not for a proposal to be presented as data.
\item \textbf{Scale.} BIZCOMPS is mostly Main Street sales; at SDE near \$1.5M (the dataset's 99th percentile is \$1.7M) and in sectors with institutional buyers (medical practices) the regression is at its edge. The extrapolation flag fires at $3\times$ beyond the fitted range, not before.
\item \textbf{Survival risk} depends on loan size, sector, state and franchise status, not on the deal's own quality; it covers SBA loans approved 2018--2026, including the period of COVID payment relief, with few loans tracked beyond about seven years. No skill score against a constant-rate baseline is reported; the cross-validated Brier score $0.0201$ should be read against that baseline (about $0.017$ for a 1.75\% event rate).
\item \textbf{SDE is an input.} Normalising SDE and judging add-backs is outside the model.
\item \textbf{Owner draw.} The financing check's DSCR uses an estimated market pay for the new owner; where the owner does most of the work a higher draw lowers the ratio.
\end{itemize}

\section*{Limits}
No number here is an appraisal or advice. The value model sees earnings, revenue and industry only; every quality adjustment moves value inside the measured spread but is not itself estimated. The survival model describes SBA-financed acquisitions, not all sales. Time and closing figures are population priors, not per-deal fits.
\end{document}
